\documentclass[11pt,a4paper]{article}

% ── Codificación, fuentes y lenguaje ───────────────────────────────────────────
\usepackage[utf8]{inputenc}
\usepackage[T1]{fontenc}
\usepackage[default,scale=0.95]{sourcesanspro}
\renewcommand{\familydefault}{\sfdefault}
\usepackage[spanish,es-noshorthands,es-tabla]{babel}

% ── Geometría y espaciado ─────────────────────────────────────────────────────
\usepackage[top=2.2cm, bottom=2.5cm, left=2.2cm, right=2.2cm, headheight=15pt]{geometry}
\usepackage{setspace}
\setstretch{1.18}
\usepackage{parskip}

% ── Matemáticas y unidades ────────────────────────────────────────────────────
\usepackage{amsmath}
\usepackage{amssymb}
\usepackage{siunitx}

% ── Circuitos ─────────────────────────────────────────────────────────────────
\usepackage[european, straightvoltages]{circuitikz}

% ── Tablas ────────────────────────────────────────────────────────────────────
\usepackage{booktabs}
\usepackage{tabularx}
\usepackage{array}
\usepackage{longtable}
\usepackage{needspace}
\usepackage{colortbl}
\usepackage{enumitem}

% ── Colores, cajas e iconos ───────────────────────────────────────────────────
\usepackage[dvipsnames]{xcolor}
\usepackage{tcolorbox}
\tcbuselibrary{skins, breakable}
\usepackage{fontawesome5}
\usepackage{graphicx}

% ── Listados de código (dark-mode infográfico) ────────────────────────────────
%   Estilo base: fondo oscuro con números de línea. Los bloques lstlisting se
%   envuelven automáticamente en una caja tcolorbox redondeada.
\usepackage{listings}

% ── Secciones con barra de color (infografía) ────────────────────────────────
\usepackage{titlesec}
\titleformat{\section}
  {\normalfont\LARGE\bfseries\color{azulNoche}}
  {\colorbox{cyanNeon}{\color{fondoOscuro}\thesection}}{0.7em}{}
  [\vspace{2pt}\color{cyanNeon}\rule{\linewidth}{1.8pt}]
\titlespacing*{\section}{0pt}{24pt}{10pt}
\titleformat{\subsection}
  {\normalfont\large\bfseries\color{azulNoche}}
  {\textcolor{cyanNeon}{\thesubsection}}{0.7em}{}
  [\vspace{1pt}\color{grisLinea}\rule{\linewidth}{0.6pt}]
\titlespacing*{\subsection}{0pt}{14pt}{6pt}
\titleformat{\subsubsection}
  {\normalfont\normalsize\bfseries\color{azulNoche}}
  {\textcolor{cyanNeon}{\thesubsubsection}}{0.7em}{}
\titlespacing*{\subsubsection}{0pt}{10pt}{4pt}

% ── Cabeceras y pies de página ────────────────────────────────────────────────
\usepackage{fancyhdr}
\usepackage{lastpage}
\pagestyle{fancy}
\fancyhf{}
\renewcommand{\headrulewidth}{0pt}
\renewcommand{\footrulewidth}{0pt}
\fancyfoot[C]{%
  \begin{minipage}{\linewidth}
    \vspace{2pt}
    \color{cyanNeon}\rule{\linewidth}{1.2pt}\\[2pt]
    \centering\color{grisTexto}\footnotesize
    Página \thepage\ de \pageref{LastPage}
  \end{minipage}%
}

% ── Hiperenlaces ──────────────────────────────────────────────────────────────
\usepackage[colorlinks=true, linkcolor=azulNoche, urlcolor=cyanNeon]{hyperref}
\sloppy
\emergencystretch 3em

% ── Paleta de colores — tecnológico dark-mode ──────────────────────────────────
\definecolor{fondoOscuro}{HTML}{0D1B2A}%             Dark navy — fondo banda título
\definecolor{verdeTurquesa}{HTML}{004D40}%           Verde turquesa — bandas de marca
\definecolor{azulNoche}{HTML}{1B3A6B}%               Midnight blue — secciones, marcos
\definecolor{azulMedio}{HTML}{2A5298}%               Azul medio — variante de acento
\definecolor{cyanNeon}{HTML}{00C8E0}%                Cyan eléctrico — acentos primarios
\definecolor{verdeSignal}{HTML}{00B887}%             Verde señal — fortalezas, éxito
\definecolor{naranjaVivo}{HTML}{FF7043}%             Naranja vivo — mejoras, advertencias
\definecolor{rojoAlerta}{HTML}{E53935}%              Rojo alerta — errores
\definecolor{grisPapel}{HTML}{F4F6FA}%               Fondo tarjetas claras
\definecolor{grisLinea}{HTML}{C8D0E0}%               Bordes sutiles
\definecolor{grisTexto}{HTML}{6B7A99}%               Texto secundario
\definecolor{amarilloNota}{HTML}{FFB300}%            Notas / advertencias doradas
\definecolor{codigoFondo}{HTML}{1E2D3D}%             Fondo del bloque de código
\definecolor{codigoComentario}{HTML}{7A8FA6}%        Comentarios
\definecolor{codigoCadena}{HTML}{FF9D5C}%            Cadenas / strings
\definecolor{codigoKeyword}{HTML}{00C8E0}%           Palabras clave
\definecolor{codigoNumero}{HTML}{00E5A0}%            Números

% ── Banda de título: portada infográfica ──────────────────────────────────────
%   Diseño en 2 capas: banda de color (por defecto fondo oscuro) + franja de
%   acento cyan abajo. Uso: \bandaTitulo[color]{título}{subtítulo}.
%   El icono se puede cambiar con \renewcommand{\iconoBanda}{...} (FontAwesome).
\newcommand{\iconoBanda}{microchip}
\newcommand{\bandaTitulo}[3][fondoOscuro]{%
  \begin{tcolorbox}[
    enhanced,
    colback=#1, colframe=#1,
    boxrule=0pt, arc=8pt, outer arc=8pt,
    width=\linewidth,
    top=18pt, bottom=6pt, left=14pt, right=14pt,
    borderline south={3.5pt}{0pt}{cyanNeon},
  ]
    \begin{minipage}[c]{0.07\linewidth}
      \centering
      \textcolor{cyanNeon}{\fontsize{30}{30}\selectfont\faIcon{\iconoBanda}}
    \end{minipage}%
    \hspace{8pt}%
    \begin{minipage}[c]{0.88\linewidth}
      {\fontsize{20}{24}\selectfont\bfseries\color{white}#2}\par\vspace{5pt}%
      \textcolor{cyanNeon!80!white}{\small\faIcon{angle-right}~#3}%
    \end{minipage}
    \vspace{4pt}
  \end{tcolorbox}%
  \vspace{8pt}%
}

% ── Tarjetas de datos (infografía) ────────────────────────────────────────────
\tcbset{
  tarjetaDato/.style={
    enhanced, breakable,
    arc=6pt, outer arc=6pt,
    colback=grisPapel, colframe=azulNoche,
    boxrule=1pt,
    fonttitle=\bfseries\small, coltitle=white,
    attach boxed title to top left={yshift=-3mm, xshift=6mm},
    boxed title style={
      colback=azulNoche, colframe=azulNoche,
      arc=4pt, boxrule=0pt,
    },
    drop fuzzy shadow=azulNoche!30!white,
    top=8pt, bottom=8pt, left=10pt, right=10pt,
  },
  tarjetaIcono/.style={
    enhanced, breakable,
    arc=6pt, outer arc=6pt,
    colback=white, colframe=grisLinea, boxrule=0.8pt,
    fonttitle=\bfseries\small, coltitle=azulNoche,
    borderline west={3pt}{0pt}{cyanNeon},
    drop fuzzy shadow=azulNoche!25!white,
    top=6pt, bottom=6pt, left=10pt, right=10pt,
  },
  cajaTitulo/.style={
    enhanced, breakable,
    arc=6pt, outer arc=6pt,
    colback=azulNoche!10!grisPapel, colframe=azulNoche,
    boxrule=1pt,
    fonttitle=\bfseries, coltitle=white,
    attach boxed title to top left={yshift=-3mm, xshift=6mm},
    boxed title style={
      colback=azulNoche, colframe=azulNoche,
      arc=4pt, boxrule=0pt,
    },
    drop fuzzy shadow=azulNoche!25!white,
    top=8pt, bottom=8pt, left=10pt, right=10pt,
  },
  cajaContenido/.style={
    enhanced, breakable,
    arc=5pt, outer arc=5pt,
    colback=grisPapel, colframe=grisLinea,
    boxrule=0.8pt,
    fonttitle=\bfseries\small, coltitle=azulNoche,
    top=6pt, bottom=6pt, left=10pt, right=10pt,
  },
  cajaFortaleza/.style={
    enhanced, breakable,
    arc=6pt, outer arc=6pt,
    colback=verdeSignal!8!white, colframe=verdeSignal,
    boxrule=1pt,
    borderline west={4pt}{0pt}{verdeSignal},
    fonttitle=\bfseries\small, coltitle=verdeSignal!70!black,
    drop fuzzy shadow=verdeSignal!20!white,
    top=6pt, bottom=6pt, left=10pt, right=10pt,
  },
  cajaMejora/.style={
    enhanced, breakable,
    arc=6pt, outer arc=6pt,
    colback=naranjaVivo!8!white, colframe=naranjaVivo,
    boxrule=1pt,
    borderline west={4pt}{0pt}{naranjaVivo},
    fonttitle=\bfseries\small, coltitle=naranjaVivo!80!black,
    drop fuzzy shadow=naranjaVivo!20!white,
    top=6pt, bottom=6pt, left=10pt, right=10pt,
  },
  cajaRecomendacion/.style={
    enhanced, breakable,
    arc=6pt, outer arc=6pt,
    colback=cyanNeon!6!white, colframe=cyanNeon!60!azulNoche,
    boxrule=1pt,
    borderline west={4pt}{0pt}{cyanNeon},
    fonttitle=\bfseries\small, coltitle=azulNoche,
    drop fuzzy shadow=azulNoche!20!white,
    top=6pt, bottom=6pt, left=10pt, right=10pt,
  },
}

% ── Ícono + texto (encabezados de sección y elementos) ────────────────────────
\newcommand{\iconotexto}[2]{\textcolor{cyanNeon}{\faIcon{#1}}~#2}

% ── Listados de código: estilo dark + auto-envoltura ─────────────────────────
\lstdefinestyle{estiloCodigo}{
    backgroundcolor=\color{codigoFondo},
    basicstyle=\ttfamily\small\color{white},
    commentstyle=\color{codigoComentario}\itshape,
    stringstyle=\color{codigoCadena},
    keywordstyle=\color{codigoKeyword}\bfseries,
    numberstyle=\tiny\color{codigoComentario},
    numbers=left,
    numbersep=10pt,
    showstringspaces=false,
    breaklines=true,
    breakatwhitespace=false,
    tabsize=4,
    frame=none,
    captionpos=b,
    aboveskip=6pt,
    belowskip=4pt,
    xleftmargin=14pt,
    framexleftmargin=14pt,
    literate=
      {á}{{\'a}}1 {é}{{\'e}}1 {í}{{\'i}}1 {ó}{{\'o}}1 {ú}{{\'u}}1
      {Á}{{\'A}}1 {É}{{\'E}}1 {Í}{{\'I}}1 {Ó}{{\'O}}1 {Ú}{{\'U}}1
      {ñ}{{\~n}}1 {Ñ}{{\~N}}1 {ü}{{\"u}}1 {Ü}{{\"U}}1
      {¿}{{?`}}1 {¡}{{!`}}1
}
\lstdefinestyle{estiloPython}{
    style=estiloCodigo,
    language=Python,
}
\lstset{style=estiloCodigo}
\BeforeBeginEnvironment{lstlisting}{%
  \begin{tcolorbox}[
    enhanced, breakable,
    arc=6pt, outer arc=6pt,
    colback=codigoFondo, colframe=azulNoche,
    boxrule=1pt,
    drop fuzzy shadow=azulNoche!25!white,
    top=2pt, bottom=2pt, left=0pt, right=0pt,
  ]%
}
\AfterEndEnvironment{lstlisting}{\end{tcolorbox}}

% ── Cajas de contenido temáticas ──────────────────────────────────────────────
\newtcolorbox{cajaRecuerda}{
    enhanced, breakable,
    arc=6pt, outer arc=6pt,
    colback=azulNoche!10!grisPapel,
    colframe=azulNoche, boxrule=1pt,
    borderline west={4pt}{0pt}{cyanNeon},
    fonttitle=\bfseries\small\color{white},
    title={\faIcon{info-circle}~Recuerda},
    attach boxed title to top left={yshift=-3mm, xshift=6mm},
    boxed title style={colback=azulNoche, arc=4pt, boxrule=0pt},
    drop fuzzy shadow=azulNoche!25!white,
    left=10pt, right=10pt, top=8pt, bottom=8pt,
}

\newtcolorbox{cajaConcepto}{
    enhanced, breakable,
    arc=6pt, outer arc=6pt,
    colback=verdeSignal!8!white,
    colframe=verdeSignal, boxrule=1pt,
    borderline west={4pt}{0pt}{verdeSignal},
    fonttitle=\bfseries\small\color{white},
    title={\faIcon{lightbulb}~Concepto clave},
    attach boxed title to top left={yshift=-3mm, xshift=6mm},
    boxed title style={colback=verdeSignal!80!black, arc=4pt, boxrule=0pt},
    drop fuzzy shadow=verdeSignal!20!white,
    left=10pt, right=10pt, top=8pt, bottom=8pt,
}

\newtcolorbox{cajaEjemplo}{
    enhanced, breakable,
    arc=6pt, outer arc=6pt,
    colback=cyanNeon!5!white,
    colframe=cyanNeon!70!azulNoche, boxrule=1pt,
    borderline west={4pt}{0pt}{cyanNeon},
    fonttitle=\bfseries\small\color{fondoOscuro},
    title={\faIcon{code}~Ejemplo},
    attach boxed title to top left={yshift=-3mm, xshift=6mm},
    boxed title style={colback=cyanNeon, arc=4pt, boxrule=0pt},
    drop fuzzy shadow=azulNoche!20!white,
    left=10pt, right=10pt, top=8pt, bottom=8pt,
}

\newtcolorbox{cajaEjercicio}{
    enhanced, breakable,
    arc=6pt, outer arc=6pt,
    colback=azulMedio!7!white,
    colframe=azulMedio, boxrule=1pt,
    borderline west={4pt}{0pt}{azulMedio},
    fonttitle=\bfseries\small\color{white},
    title={\faIcon{pencil-alt}~Ejercicio},
    attach boxed title to top left={yshift=-3mm, xshift=6mm},
    boxed title style={colback=azulMedio, arc=4pt, boxrule=0pt},
    drop fuzzy shadow=azulNoche!20!white,
    left=10pt, right=10pt, top=8pt, bottom=8pt,
}

\newtcolorbox{cajaReto}{
    enhanced, breakable,
    arc=6pt, outer arc=6pt,
    colback=naranjaVivo!6!white,
    colframe=naranjaVivo, boxrule=1pt,
    borderline west={4pt}{0pt}{naranjaVivo},
    fonttitle=\bfseries\small\color{white},
    title={\faIcon{fire}~Reto},
    attach boxed title to top left={yshift=-3mm, xshift=6mm},
    boxed title style={colback=naranjaVivo, arc=4pt, boxrule=0pt},
    drop fuzzy shadow=naranjaVivo!20!white,
    left=10pt, right=10pt, top=8pt, bottom=8pt,
}

% ── Indicadores de dificultad ─────────────────────────────────────────────────
\newcommand{\facil}{\textcolor{verdeSignal}{\faIcon{circle}\,\textbf{Fácil}}}
\newcommand{\intermedio}{\textcolor{naranjaVivo}{\faIcon{circle}\,\textbf{Intermedio}}}
\newcommand{\dificil}{\textcolor{rojoAlerta}{\faIcon{circle}\,\textbf{Difícil}}}

% ─────────────────────────────────────────────────────────────────────────────
\renewcommand{\iconoBanda}{shield-alt}
\fancyhead[L]{\footnotesize\color{grisTexto}\faIcon{shield-alt}~CYBERSEGURIDAD}
\fancyhead[R]{\footnotesize\color{grisTexto}Semana 02}
\hypersetup{pdftitle={Ciberseguridad Aplicada},pdfauthor={Prof. César Rodríguez}
  pdfsubject={CYBERSEGURIDAD}}

\begin{document}
\bandaTitulo[fondoOscuro]{Semana 02 --- Criptografía aplicada y gestión de secretos}{Fase 1: Fundamentos, criptografía y marco legal --- instrucciones paso a paso}

\begin{center}\small\color{grisTexto}Prof. César Rodríguez \textbar{} ELECTRONICA \textbar{} 1 de noviembre de 2026\end{center}

\begin{tabularx}{\linewidth}{@{}>{\raggedright\arraybackslash}p{4.6cm}X@{}}
\toprule
\textbf{Curso} & CYBERSEGURIDAD --- Ciberseguridad Aplicada \\
\textbf{Semana} & 02 de 10 \\
\textbf{Fase} & 1 --- Fundamentos, criptografía y marco legal \\
\textbf{Modalidad} & 4 grupos, trabajo desde sus propios equipos \\
\textbf{Dedicación} & 4 horas de laboratorio + 2 horas de estudio independiente \\
\textbf{Instrumento} & Actividad formativa \\
\textbf{Entrega} & GitHub, carpeta semana\_02/ \\
\textbf{Corte} & Domingo 23:59 de la semana correspondiente; lo posterior no se corrige. \\
\bottomrule
\end{tabularx}

\section{\iconotexto{clock}{Antes de la sesión}}

Estas tareas se hacen en casa, antes de la práctica. Si el repositorio o el laboratorio no están listos al empezar, la sesión se va en arreglarlos y no alcanza para el informe.

\begin{enumerate}[leftmargin=1.4em,itemsep=2pt,topsep=3pt]
  \item Comprobar que el repositorio del grupo tiene la carpeta semana\_02 creada.
  \item Revisar la retroalimentación del docente sobre la semana anterior y corregir lo pendiente.
  \item Traer instaladas las herramientas de la semana: openssl, sha256sum, Python hashlib.
  \item Confirmar que el laboratorio arranca y que los tres servicios responden antes de la sesión.
  \item Leer el objetivo y los contenidos para llegar con el vocabulario de la práctica.
\end{enumerate}

\section{\iconotexto{bullseye}{Objetivo de la semana}}

Aplicar cifrado simétrico y asimétrico, funciones hash y derivación de claves, y diagnosticar una configuración TLS insegura.

\section{\iconotexto{clipboard-list}{Contenidos que se aplican}}

\begin{itemize}[leftmargin=1.4em,itemsep=2pt,topsep=3pt]
  \item Cifrado simétrico AES en modos GCM y CBC, y por qué GCM aporta integridad.
  \item Cifrado asimétrico RSA y curvas elípticas; intercambio de claves.
  \item Hash criptográfico, HMAC y funciones de derivación PBKDF2 y Argon2.
  \item TLS 1.2 y 1.3, cadena de certificados y errores de validación.
  \item Gestión de secretos y detección de claves en repositorios.
\end{itemize}

\section{\iconotexto{flask}{Escenario}}

Cada grupo cifra y descifra un archivo, firma y verifica un mensaje, inspecciona un certificado real y mide el coste de derivar una clave con distinto número de iteraciones.


\section{\iconotexto{list-ol}{Paso a paso de la práctica}}

\begin{enumerate}[leftmargin=1.4em,itemsep=2pt,topsep=3pt]
  \item Cifrar y descifrar un archivo con AES-256-GCM y verificar que un byte alterado rompe la integridad.
  \item Generar un par de claves, firmar un mensaje y verificar la firma; después alterar el mensaje y comprobar el fallo.
  \item Extraer e inspeccionar el certificado de un sitio real, identificando emisor, validez y algoritmo.
  \item Medir el tiempo de PBKDF2 con 10000, 100000 y 600000 iteraciones y tabular el resultado.
  \item Documentar qué primitiva resuelve confidencialidad, integridad y autenticidad en cada caso.
\end{enumerate}

\begin{cajaConcepto}\textbf{Regla de oro de la evidencia:} guarda la salida de cada paso en \texttt{semana\_02/evidencia/} en el momento en que la obtienes. Reconstruirla después no es evidencia, es memoria.
\end{cajaConcepto}

\textbf{Herramientas de la semana:} openssl, sha256sum, Python hashlib.

\section{\iconotexto{users}{Rol de cada grupo esta semana}}

Todos los grupos resuelven la misma práctica, cada uno en su propio repositorio y su propio laboratorio: el trabajo de un grupo no depende del de otro. El rol cambia el énfasis del trabajo y lo que cada grupo aporta al comparar resultados en clase, al final.

\begin{tabularx}{\linewidth}{@{}l l X@{}}
\toprule
Grupo & Rol & Aporte principal\\
\midrule
Grupo 1 & \textbf{Rojo} & Evidencia cruda en semana\_02/evidencia y el apartado de reproducción del informe. \\
Grupo 2 & \textbf{Azul} & La mitigación aplicada, el comando de verificación y la comparación de resultados. \\
Grupo 3 & \textbf{Púrpura} & La regla o consulta de detección y la tabla de lo detectado frente a lo ejecutado. \\
Grupo 4 & \textbf{Legal y reporte} & La versión final del informe y la declaración explícita de lo que no se pudo verificar. \\
\bottomrule
\end{tabularx}

\section{\iconotexto{box-open}{Qué se entrega}}

Registro de comandos, salidas y una tabla que asigne cada primitiva a la propiedad que garantiza.

\textbf{Estructura de la carpeta} \texttt{semana\_02/}\textbf{:}

\begin{itemize}[leftmargin=1.4em,itemsep=2pt,topsep=3pt]
  \item README.md con integrantes, roles y tabla de semanas entregadas.
  \item semana\_02/informe.md o informe.pdf con el entregable de esa práctica.
  \item semana\_02/scripts/ con el código, con type hints y docstrings.
  \item semana\_02/evidencia/ con salidas, capturas y hashes anonimizados.
  \item semana\_02/README.md con el resumen de lo hecho y lo que no se pudo verificar.
\end{itemize}

\textbf{Reglas que se verifican al corregir:}

\begin{itemize}[leftmargin=1.4em,itemsep=2pt,topsep=3pt]
  \item Nada de secretos, tokens ni contraseñas en el repositorio; las claves van en .env ignorado.
  \item Nada de imágenes de disco, muestras de malware ni binarios: solo hashes y enlaces.
  \item Evidencia anonimizada: sin datos personales de terceros ni credenciales reales.
  \item Un commit por entrega con mensaje "semana NN: \textless{}entregable\textgreater{}".
  \item La práctica cuenta como entregada solo si está en el repositorio antes del corte.
\end{itemize}

\section{\iconotexto{ruler-combined}{Cómo se corrige}}

\begin{tabularx}{\linewidth}{@{}X r@{}}
\toprule
Criterio & Peso\\
\midrule
Cifrado y descifrado correctos con verificación de integridad & 3,0 \\
Firma y verificación demostradas, incluido el caso alterado & 2,5 \\
Lectura correcta del certificado y de la cadena & 2,0 \\
Medición de KDF con conclusión sobre coste y seguridad & 1,5 \\
Explicación de qué propiedad cubre cada primitiva & 1,0 \\
\midrule
\textbf{Total} & \textbf{10,0}\\
\bottomrule
\end{tabularx}

\section{\iconotexto{clipboard-check}{Lista de verificación antes del corte}}

\begin{itemize}[leftmargin=1.4em,itemsep=2pt,topsep=3pt]
  \item El repositorio del grupo existe, es privado y el docente aparece como colaborador.
  \item La carpeta semana\_02/ está en el repositorio, no solo en el equipo.
  \item El informe explica qué se hizo, con qué comandos y qué se obtuvo.
  \item La evidencia está anonimizada: sin datos personales, credenciales ni secretos.
  \item Cada hallazgo tiene su hash o su salida cruda que permite reproducirlo.
  \item El commit 'semana 02: \textless{}entregable\textgreater{}' está hecho antes del corte.
\end{itemize}

\begin{tcolorbox}[cajaMejora]\textbf{Alcance y autorización:} solo se trabaja sobre los contenedores del laboratorio o sobre sistemas con autorización escrita. Si el objetivo no está en ese conjunto, se detiene la práctica y se consulta al docente. El uso de estas técnicas fuera del alcance tiene consecuencias académicas y legales.
\end{tcolorbox}

\begin{cajaRecuerda}Dudas de la semana: se responden en clase o en el repositorio. Corte: Domingo 23:59 de la semana correspondiente; lo posterior no se corrige.
\end{cajaRecuerda}

\clearpage

\section{\iconotexto{book}{Glosario de la semana}}

\small
\begin{longtable}{@{}>{\raggedright\arraybackslash}p{4.2cm}>{\raggedright\arraybackslash}p{\dimexpr\linewidth-4.2cm-12pt\relax}@{}}
\toprule
\textbf{Término} & \textbf{Definición}\\
\midrule
\endhead
\bottomrule
\endlastfoot
\textbf{Cifrado simétrico} & Cifrado que usa la misma clave para cifrar y descifrar. Es rápido y adecuado para volumen, pero exige repartir la clave por un canal seguro.\\
\textbf{Cifrado asimétrico} & Cifrado con un par de claves: la pública cifra o verifica y la privada descifra o firma. Resuelve el reparto de claves a cambio de ser más lento.\\
\textbf{Hash criptográfico} & Función de un solo sentido que produce una huella de tamaño fijo; cambiar un byte cambia la huella. No cifra ni se puede revertir.\\
\textbf{HMAC} & Hash con clave que autentica un mensaje y detecta manipulación. Aporta integridad y autenticidad de origen, no confidencialidad.\\
\textbf{Derivación de claves} & Convertir una contraseña o secreto débil en una clave del tamaño adecuado con un coste deliberado (PBKDF2, Argon2) para frenar la fuerza bruta.\\
\textbf{TLS} & Protocolo que cifra y autentica una conexión entre cliente y servidor. HTTPS es HTTP sobre TLS.\\
\textbf{Cadena de certificados} & Secuencia de certificados firmados que enlaza la clave del servidor con una autoridad de confianza. Si falta un eslabón o es inválido, la validación falla.\\
\textbf{Gestión de secretos} & Crear, almacenar, rotar y revocar claves y contraseñas sin dejarlas en el código, en los argumentos ni en el historial de git.\\
\end{longtable}
\end{document}
